%=============================================================================!
% 4.3  THE PHASE PORTRAIT                                                     !
%=============================================================================!
\section{The phase portrait}
\label{sec:os-phase}

\Cref{sec:ne-equations} showed that the state of a body moving along a
line is the pair $(x, v)$, and drew motions as paths through the plane of
states (\cref{fig:ne-state}b). This plane is called the phase plane, and a
drawing of the paths through it the phase portrait of the system. For the
oscillator, energy alone gives the shape of every path.

\subsection*{Ellipses of constant energy}

Along any motion of the block, $\frac12 mv^2 + \frac12 kx^2 = E$. Dividing
both sides by $E$, which is positive unless the block rests at $x = 0$,
gives $kx^2/2E + mv^2/2E = 1$, and dividing the top and bottom of the
first fraction by $k$ and of the second by $m$,
\begin{equation}
   \frac{x^2}{2E/k} + \frac{v^2}{2E/m} = 1 ,
   \qquad\text{or}\qquad
   \Bigl(\frac{x}{a}\Bigr)^2 + \Bigl(\frac{v}{b}\Bigr)^2 = 1
   \quad\text{with } a = \sqrt{\frac{2E}{k}},\ b = \sqrt{\frac{2E}{m}} .
   \label{eq:os-ellipse}
\end{equation}
The second form says that the point $(x/a, v/b)$ lies on the circle of
radius 1, so the state $(x, v)$ lies on that circle stretched by the
factor $a$ along the $x$ axis and $b$ along the $v$ axis. A stretched
circle is an ellipse, and $a$ and $b$ are its semi-axes, half its width
and half its height. Here $a = \sqrt{2E/k}$ is the amplitude $A$
(\cref{sec:os-energy}), and $b = \sqrt{2E/m} = \sqrt{(2E/k)(k/m)} =
A\omega$, multiplying and dividing by $k$ under the root, is the largest
speed. For the block with $E =
\SI{0.04}{\joule}$ they are \SI{4}{\centi\metre} and
\SI{0.4}{\metre\per\second}.

\Cref{fig:os-phase}(a) draws the ellipses for three energies. Every
state lies on exactly one such ellipse, the one of its own energy, and the
block goes round it for ever, because energy is conserved. Two paths can
never cross: if they met at one state, that state would start two
different motions, which \cref{sec:ne-equations} rules out. The
direction of travel is clockwise: in the upper half of the plane $v > 0$,
so $x$ grows and the state moves to the right; in the lower half it moves
to the left. The panel is drawn with \SI{1}{\centi\metre} of position as
long as \SI{0.1}{\metre\per\second} of velocity, that is, with the velocity
divided by $\omega = \SI{10}{\radian\per\second}$, since
$\SI{0.1}{\metre\per\second}/\SI{10}{\radian\per\second} =
\SI{0.01}{\metre}$. On that scale the ellipses are circles, since $(x,
v/\omega) = A(\cos\theta, -\sin\theta)$ by \cref{sec:os-circle}: the
reference circle again, reflected in the $x$ axis, which reverses the sign
of the second coordinate and so turns the anticlockwise motion of $P$ into
a clockwise one.

\begin{figure}[htb]
\centering
\includegraphics{ch04_oscillations/phase}
\caption{Phase portraits. (a) The block of \cref{sec:os-circle} at the
energies 0.01, 0.04 (teal) and \SI{0.09}{\joule}, traversed clockwise;
with \SI{1}{\centi\metre} drawn as long as \SI{0.1}{\metre\per\second}
the ellipses are circles. (b) The pendulum of \cref{sec:os-anharmonic},
with the rate of change of its angle, $\dot\theta$, in units of its
small-swing angular frequency $\omega_0$: closed paths while it swings (grey), the path through the
upright position (teal), and open paths when it goes over the top
(ochre).}
\label{fig:os-phase}
\end{figure}

Phase portraits return throughout the book. Chapter~7 uses momentum in
place of velocity and studies the flow of whole regions of the phase
plane, and Chapter~9 judges methods of computing motion by whether their
paths close, as these do, or spiral outwards. Panel (b), the pendulum,
belongs to \cref{sec:os-anharmonic}; there $\theta$ is the angle of the
pendulum, not $\omega t + \phi$. Its teal paths seem to cross at the
upright position, but that state is a path of its own, a pendulum
balanced at rest, which a pendulum on the teal paths approaches ever more
slowly and never reaches, so no state starts two motions.

\begin{selfcheck}
A block with $\omega = \SI{10}{\radian\per\second}$ starts at $x = 0$
with $v = \SI{0.3}{\metre\per\second}$. What are the semi-axes of its
ellipse in the phase plane, and where on the ellipse does it start?
\end{selfcheck}
